\documentclass{article}

\usepackage[english]{babel}
\usepackage[letterpaper,top=2.5cm,bottom=2.5cm,left=3cm,right=3cm]{geometry}
\usepackage{amsmath}
\usepackage{graphicx}
\usepackage{booktabs}
\usepackage[colorlinks=true, allcolors=blue]{hyperref}

\title{Summary Statistics: Koijen \& Yogo (2020) Replication Data}
\author{Replication Project}
\date{}

\begin{document}
\maketitle

\begin{abstract}
This document summarises the underlying data used in the replication of
Table~1 of Koijen \& Yogo (2020), ``Exchange Rates and Asset Prices in a
Global Demand System.''
The exhibits below provide a cross-sectional snapshot of outstanding financial
assets for the latest year available, a time-series view of total amounts
by asset class, a country-level breakdown of total outstanding, and a
cross-country view of foreign ownership shares.  Together they allow the
reader to assess data coverage, the relative scale of each asset class,
which economies dominate the sample, and how internationally held each
market's assets are.
\end{abstract}

\section{Summary Statistics Table}

Table~\ref{tab:summary_stats} presents cross-sectional summary statistics for
outstanding financial-asset amounts across the 37-country sample used in the
replication.
Data are sourced from OECD Table 720 (debt and equity for OECD members),
BIS Domestic Debt Securities Statistics (non-OECD fallback for debt), and
World Bank WDI market-capitalisation data (non-OECD fallback for equity).

The key takeaway is that long-term debt is by far the largest asset class
by aggregate outstanding, reflecting the importance of sovereign and
corporate bond markets globally.  Equity exhibits the widest cross-country
dispersion (high standard deviation relative to mean), driven by the
outsized size of US, Japanese, and UK equity markets versus smaller
emerging-market economies.

\input{../_output/summary_stats_table.tex}

\section{Time-Series Chart}

Figure~\ref{fig:summary_chart} plots the total outstanding amounts for each
asset class---short-term debt, long-term debt, and equity---summed across all
countries in the 37-country sample, expressed in USD trillions.

\begin{figure}[ht]
\centering
\includegraphics[width=0.92\textwidth]{../_output/summary_stats_chart.png}
\caption{%
  \textbf{Total outstanding financial assets by asset class, 37-country sample
  (Koijen \& Yogo 2020 replication).}
  The figure shows USD-trillion totals for short-term debt (type~1, blue),
  long-term debt (type~2, green), and equity (type~3, pink) from the earliest
  available year through the latest API vintage.
  The reader should take away three facts:
  (i)~long-term debt dominates total portfolio size throughout the sample;
  (ii)~equity grew steeply in the 2010s before the post-2021 correction,
  reflecting global equity-market cycles;
  (iii)~short-term debt is small relative to the other two classes, consistent
  with the paper's finding that foreign investor demand is concentrated in
  long-term instruments.
  Data sources: OECD T720, BIS WS\_NA\_SEC\_DSS, World Bank WDI.
}
\label{fig:summary_chart}
\end{figure}

\section{Country-Level Outstanding Amounts}

Figure~\ref{fig:country_bar} ranks all 37 sample countries by their total
outstanding financial assets in the latest available year, stacked by asset class.

\begin{figure}[ht]
\centering
\includegraphics[width=0.92\textwidth]{../_output/summary_stats_country_bar.png}
\caption{%
  \textbf{Outstanding financial assets by country and asset class, latest year
  (Koijen \& Yogo 2020 replication).}
  Each bar shows a country's total outstanding (USD trillions) split into
  short-term debt (blue), long-term debt (green), and equity (pink).
  Countries are ranked from smallest to largest total.
  The reader should take away that the United States dwarfs all other
  economies in both long-term debt and equity, followed by Japan and the
  euro-area economies.  Equity is absent or very small for countries where
  OECD T720 F5 data are unavailable and only World Bank market-cap data
  are used, exposing a structural coverage gap for smaller emerging markets.
  Data sources: OECD T720, BIS WS\_NA\_SEC\_DSS, World Bank WDI.
}
\label{fig:country_bar}
\end{figure}

\section{Foreign Ownership Shares}

Figure~\ref{fig:foreign_share} shows the foreign ownership share --- defined
as total cross-border holdings by all foreign investors divided by total
outstanding --- for each country and asset class in the latest year.
This is the key quantity studied by Koijen \& Yogo (2020): understanding
\emph{who} holds which country's assets and \emph{how much} is held abroad
is the foundation of their global demand-system model.

\begin{figure}[ht]
\centering
\includegraphics[width=0.92\textwidth]{../_output/summary_stats_foreign_share.png}
\caption{%
  \textbf{Foreign ownership share by country and asset class, latest year
  (Koijen \& Yogo 2020 replication).}
  Each dot shows the share of outstanding assets of a given class held by
  foreign investors, as measured by IMF PIP/CPIS bilateral position data.
  Circles = short-term debt, squares = long-term debt, triangles = equity.
  Countries are ranked by their average foreign share across asset classes.
  The reader should take away that small open economies (e.g.\ Belgium,
  the Netherlands, Singapore) have the highest foreign ownership shares,
  consistent with the paper's finding that these countries face the most
  elastic foreign demand.  The United States and Japan, despite their
  enormous market size, have relatively low foreign shares, especially in
  equity, reflecting home-bias in global portfolios.
  Foreign shares computed from IMF CPIS; outstanding from OECD/BIS/World Bank.
}
\label{fig:foreign_share}
\end{figure}

\section{Data Coverage Notes}

\begin{itemize}
  \item \textbf{Debt}: OECD T720 (DF\_T720R\_A) is the primary source for
    27 OECD members; BIS WS\_NA\_SEC\_DSS fills non-OECD economies
    (China, India, Malaysia, Philippines, Russia, South Africa, Thailand,
    Hong Kong, Singapore, Australia).
  \item \textbf{Equity}: OECD T720 instrument F5 (equity + investment fund
    shares) is used for OECD members; World Bank market-capitalisation
    data (CM.MKT.LCAP.CD) fills non-OECD.  The OECD F5 measure is known
    to overstate market capitalisation by 2--3$\times$ because it includes
    investment fund shares; the paper uses proprietary Datastream data which
    is unavailable in this free-data replication.
  \item \textbf{Bilateral holdings}: IMF Portfolio Investment Position (PIP/CPIS)
    bilateral parquets provide the foreign-ownership share denominators used
    in Tables 1 and 2.
  \item \textbf{Known gaps}: New Zealand debt is absent from all free sources;
    China only reports long-term general-government debt via BIS.
\end{itemize}

\end{document}
